%% ============================================================
%% it/sections/08_passione.tex
%% Layer opzionale: Pasión de las Pasiones
%% ============================================================

\chapter{La Passione dei Vespri}
\markboth{La Passione dei Vespri}{La Passione dei Vespri}

\lettrine[lines=3]{L}{a} storia dei Vespri Siciliani è, tra le altre cose,
una storia di passioni. Amori impossibili tra famiglie nemiche. Tradimenti
che nascono dall'amore e finiscono nel sangue. Genitori che sacrificano
i figli alla causa. Rivali che si odiano con la stessa intensità con cui
una volta si amavano.

Questo capitolo presenta un \textbf{layer drammatico opzionale} ispirato
a \textit{Pasión de las Pasiones} di Brandon León-Gambetta. Non sostituisce
il sistema FitD — lo affianca. Ogni personaggio può avere, in aggiunta
al proprio playbook, un \textbf{Ruolo Passionale}: un secondo strato
identitario che definisce la natura del suo dramma personale e garantisce
che, prima o poi, quel dramma esploda.

\nota[Per il GM — Quando usare questo layer]{
  Il layer Passionale è opzionale e si attiva tavolo per tavolo. Funziona
  meglio con giocatori che amano il dramma personale oltre all'azione
  politica. Se il tuo tavolo è più orientato alla strategia che alle
  relazioni, puoi ignorarlo completamente senza perdere nulla di essenziale.
  Se invece vuoi che i tuoi giocatori piangano — e ridano — attivalo.
}

\secsep

\section{Come Funziona}

\subsection{Il Ruolo Passionale}

All'inizio della sessione, ogni giocatore che vuole partecipare al layer
sceglie un Ruolo Passionale dalla lista in questo capitolo. Il Ruolo è
indipendente dal playbook: una Castellana può essere L'Amore Impossibile;
un Frate Predicatore può essere Il Traditore.

Ogni Ruolo ha:
\begin{itemize}[nosep]
  \item \textbf{Il Tormento} — la tensione emotiva che definisce il personaggio
  \item \textbf{Il Trigger} — la situazione che fa esplodere il dramma
  \item \textbf{La Promessa} — qualcosa che \textit{deve} accadere prima
    della fine della sessione (il GM lo garantisce narrativamente)
  \item \textbf{Il Momento} — una scena speciale che il giocatore può
    richiedere una volta sola
  \item \textbf{Il Punteggio Passione} — una risorsa separata dallo Stress
\end{itemize}

\subsection{Il Punteggio Passione}

Ogni personaggio con un Ruolo Passionale ha \textbf{5 caselle Passione}.
La Passione è distinta dallo Stress: rappresenta l'intensità emotiva
del momento, non il danno fisico o psicologico.

\regola[Guadagnare Passione]{
  Guadagni 1 Passione quando:
  \begin{itemize}[nosep]
    \item il tuo Trigger si attiva
    \item qualcuno che ami è in pericolo per colpa tua
    \item devi scegliere tra la causa e il tuo cuore
    \item qualcuno scopre un tuo segreto
  \end{itemize}
}

\regola[Spendere Passione]{
  Puoi spendere Passione (1 casella) per:
  \begin{itemize}[nosep]
    \item aggiungere +1d a qualsiasi tiro che coinvolga emozioni,
      relazioni o protezione di qualcuno amato
    \item dichiarare un \textbf{Momento Drammatico}: il tuo personaggio
      fa o dice qualcosa di memorabile — il GM deve trovare spazio
      narrativo per questo momento
    \item rifiutare le conseguenze di un tiro che riguarda il tuo Tormento,
      sostituendole con una conseguenza emotiva invece che fisica
  \end{itemize}
}

\regola[Passione al Massimo]{
  Quando tutte e 5 le caselle sono piene, il personaggio entra in
  \textbf{Crisi Passionale}: deve agire secondo il proprio Tormento,
  anche se è controproducente. Dopo la Crisi, le caselle si svuotano.
  Una Crisi Passionale ben giocata vale automaticamente +1 al punteggio
  della fazione — perché il dramma personale, in questo gioco, muove
  anche la storia.
}

\secsep

\section{I Ruoli Passionali}

\subsection{La Protagonista}
\textit{Qualcuno ha deciso il tuo destino. Tu hai deciso di no.}

\textbf{Tormento:} Sei stata/o definita/o da ciò che altri volevano da te
— il tuo nome, la tua famiglia, il tuo Dio, la tua causa. Stai scoprendo,
proprio adesso, in mezzo a questa crisi, che forse sei qualcosa di diverso.
Il problema è che non sai ancora cosa.

\textbf{Trigger:} Qualcuno ti dice chi sei o cosa devi fare. Ogni volta
che questo accade, guadagni 1 Passione.

\textbf{La Promessa:} Prima della fine della sessione, farai una scelta
che contraddice ciò che tutti si aspettano da te. Non sarà la scelta
giusta. Sarà la tua.

\textbf{Il Momento:} Una volta sola, puoi dichiarare: \textit{«Adesso
basta.»} Il tuo personaggio si ferma, si volta, e dice o fa qualcosa che
nessuno si aspettava. Tiri 0 dadi — il risultato è automaticamente un
successo pieno, ma la conseguenza narrativa è irreversibile.

\textbf{Adatta al 1282:} La Moglie del Traditore che decide di non
nascondersi più. Il Figlio che disobbedisce al padre nel momento
cruciale. La Badessa che apre le porte del convento sapendo cosa
ne seguirà.

\filetto

\subsection{L'Amore Impossibile}
\textit{Lo ami. Non puoi amarlo. Continui ad amarlo.}

\textbf{Tormento:} Ami qualcuno che non dovresti — per classe, per fazione,
per fede, per il sangue versato tra le vostre famiglie. Non è una
scelta. Non è una debolezza. È semplicemente vero, e il mondo non
lascia spazio a questa verità.

\textbf{Trigger:} Ogni volta che quella persona è in scena con te,
guadagni 1 Passione. Ogni volta che la metti in pericolo con le
tue azioni, guadagni 2 Passione.

\textbf{La Promessa:} Prima della fine della sessione, avrete un
momento da soli — anche se breve, anche se interrotto. Il GM lo
garantisce. Ciò che succede in quel momento dipende da voi.

\textbf{Il Momento:} Una volta sola, puoi attraversare qualsiasi
linea — fazione, pericolo, senso comune — per proteggere o raggiungere
quella persona. Tiri normalmente, ma qualunque sia il risultato,
ci arrivi. Il costo lo paghi dopo.

\textbf{Adatta al 1282:} Il Cavaliere angioino che ama una donna
siciliana. Il Frate che ama qualcuno che non può avere per ragioni
di voto. Il Congiurato innamorato di qualcuno nella fazione sbagliata.

\nota[Nota cross-tavolo]{
  L'Amore Impossibile funziona meglio se la persona amata è un PG
  di un altro tavolo. Il GM coordinatore può gestire i momenti
  condivisi come eventi speciali tra i Momenti di Congiunzione.
  Se non è possibile logisticamente, la persona amata può essere
  un PNG significativo.
}

\filetto

\subsection{Il Traditore}
\textit{Hai un segreto. Il segreto distruggerà tutto. Non riesci a smettere di portarlo.}

\textbf{Tormento:} Stai facendo qualcosa — o hai fatto qualcosa — che
tradirebbe tutto ciò che gli altri si aspettano da te. Non sei
necessariamente un villain: forse lo fai per amore, per paura,
per un debito impossibile da pagare. Ma se viene fuori, nulla
sarà più come prima.

\textbf{Trigger:} Ogni volta che qualcuno si avvicina alla verità,
guadagni 1 Passione. Ogni volta che menti attivamente su ciò che
stai facendo, guadagni 1 Passione.

\textbf{La Promessa:} Il segreto verrà fuori prima della fine della
sessione. Non necessariamente davanti a tutti — ma almeno davanti
a una persona che non avrebbe dovuto sapere. Il GM lo garantisce.
Il come dipende dalle tue azioni.

\textbf{Il Momento:} Una volta sola, puoi scegliere di confessare
volontariamente il tuo segreto a qualcuno. Se lo fai, svuoti tutte
le caselle Passione e guadagni +2 dadi al prossimo tiro — qualunque
esso sia. La confessione non cambia i fatti. Cambia chi sei.

\textbf{Adatta al 1282:} La Spia che fa il doppio gioco anche con
i Congiurati. Il Vassallo che informa i francesi per proteggere
la famiglia. Il Canonico che conosce qualcosa di devastante e
non lo ha mai detto.

\filetto

\subsection{La Figura Paterna/Materna}
\textit{Hai dato tutto per proteggere qualcuno. Stanotte potrebbe non bastare.}

\textbf{Tormento:} Qualcuno dipende da te — un figlio, un nipote,
un'intera comunità, un giovane che hai allevato come tuo. Hai
costruito tutta la tua vita intorno a questa protezione. E adesso
la situazione ti chiede di scegliere tra loro e la causa.

\textbf{Trigger:} Ogni volta che la persona che proteggi è in
pericolo, guadagni 2 Passione. Ogni volta che devi scegliere tra
lei e la causa, guadagni 1 Passione.

\textbf{La Promessa:} Prima della fine della sessione, la persona
che proteggi sarà in pericolo reale — non per colpa tua, ma a causa
degli eventi. Come la salvi (o non la salvi) è il cuore della
tua storia.

\textbf{Il Momento:} Una volta sola, puoi interporre te stesso/a
tra la persona amata e il pericolo. Assorbi automaticamente tutto
il danno che avrebbe subito. Nessun tiro. Poi tiri Resistere per
vedere quanto ti è costato.

\textbf{Adatta al 1282:} La Pescivendola e il Ragazzo. Il Vecchio
e chiunque gli ricordi qualcuno. Il Vescovo e il suo gregge. Il Barone
e il Figlio che vuole combattere.

\filetto

\subsection{Il Rivale}
\textit{Vuole ciò che hai tu. O forse è il contrario. A questo punto non lo ricordate più.}

\textbf{Tormento:} C'è qualcuno — un altro PG o un PNG significativo
— con cui hai una storia. Forse eravate amici, amanti, fratelli.
Ora siete avversari. Non odio puro: qualcosa di più complicato,
dove il rispetto e il risentimento si mescolano in proporzioni
che cambiano ogni scena.

\textbf{Trigger:} Ogni volta che il tuo Rivale ottiene qualcosa che
tu volevi, guadagni 1 Passione. Ogni volta che lo superi, guadagni
1 Passione (sì, la vittoria costa quanto la sconfitta).

\textbf{La Promessa:} Prima della fine della sessione, avrete un
confronto diretto — non necessariamente violento. Qualcosa verrà
detto che non poteva restare non detto.

\textbf{Il Momento:} Una volta sola, puoi dichiarare di aiutare
il tuo Rivale invece di ostacolarlo. Questo non cambia la vostra
relazione — la complica. Guadagni 3 Passione e il Rivale ti deve
qualcosa che non può rifiutare.

\textbf{Adatta al 1282:} Due baroni con un vecchio conto aperto.
Il Cospiratore e la Spia che una volta lavoravano insieme. Il Frate
Predicatore e un predicatore della parte opposta.

\filetto

\subsection{L'Innocente Coinvolto}
\textit{Non volevi essere qui. Adesso non puoi più andartene.}

\textbf{Tormento:} Sei finito/a in questa storia per caso, per
sfortuna, o perché qualcuno che amavi ti ha trascinato. Non hai
grandi ideali. Non hai un piano. Hai solo la sensazione crescente
che stanotte cambierà tutto — e che tu non sarai più la stessa
persona dopo.

\textbf{Trigger:} Ogni volta che vedi qualcosa che non avresti
dovuto vedere, guadagni 1 Passione. Ogni volta che qualcuno che
credevi buono fa qualcosa di terribile, guadagni 1 Passione.

\textbf{La Promessa:} Prima della fine della sessione, farai
qualcosa di coraggioso. Non perché sei un eroe — ma perché non
potevi fare altrimenti. Il GM garantisce che si presenti il momento.

\textbf{Il Momento:} Una volta sola, puoi dichiarare di esserti
trovata/o nel posto giusto al momento giusto. Qualcosa che stava
per andare molto storto non va storto — perché tu eri lì. Nessun
tiro. Il GM integra narrativamente.

\textbf{Adatta al 1282:} Il Ragazzo che non capisce ancora cosa
sta succedendo. Il Pellegrino che avrebbe dovuto ripartire ieri.
Il Medico Arabo che voleva solo curare qualcuno.

\filetto

\subsection{La Martire}
\textit{Sai già come va a finire. Vai avanti lo stesso.}

\textbf{Tormento:} Hai visto abbastanza da capire che la tua strada
porta verso qualcosa di doloroso — forse la morte, forse la perdita
di tutto ciò che ami, forse semplicemente la fine della persona che
eri. Non stai cercando di evitarlo. Stai cercando di farlo valere.

\textbf{Trigger:} Ogni volta che qualcuno cerca di fermarti o
proteggerti, guadagni 1 Passione. Ogni volta che scegli il percorso
più duro quando ce n'è uno più sicuro, guadagni 1 Passione.

\textbf{La Promessa:} Prima della fine della sessione, ti verrà
offerta una via d'uscita. Potrai scegliere se prenderla o no.

\textbf{Il Momento:} Una volta sola, puoi fare qualcosa di
irreversibile che garantisce il successo della causa — a un costo
personale che definisci tu. Il GM non può negare l'effetto. Il costo
non può essere negato da nessuno, nemmeno da te.

\textbf{Adatta al 1282:} Il Nobile in Esilio che sa che tornerà
in patria solo per morire. Il Frate Predicatore che predica sapendo
che lo arresteranno. La Vedova che non ha più niente da perdere.

\filetto

\subsection{Il Redento}
\textit{Hai fatto cose di cui ti vergogni. Stanotte potresti fare qualcosa di diverso.}

\textbf{Tormento:} Il tuo passato ti pesa. Forse hai collaborato
con gli Angioini, forse hai tradito qualcuno, forse hai semplicemente
guardato dall'altra parte troppe volte. Non stai cercando il perdono
degli altri — stai cercando di capire se sei ancora capace di
meritarlo.

\textbf{Trigger:} Ogni volta che qualcuno ti tratta come se fossi
ancora quella persona che eri, guadagni 1 Passione. Ogni volta che
hai l'occasione di fare la cosa giusta e la fai davvero, guadagni
1 Passione.

\textbf{La Promessa:} Prima della fine della sessione, il tuo
passato tornerà — qualcuno, qualcosa, una conseguenza di ciò che
hai fatto. Come lo affronti definisce chi sei adesso.

\textbf{Il Momento:} Una volta sola, puoi dichiarare di fare
qualcosa che il tuo vecchio sé non avrebbe mai fatto. Tiri con
vantaggio (+1d) e il GM deve riconoscere narrativamente la differenza.

\textbf{Adatta al 1282:} Il Soldato Disertore. La Moglie del
Traditore. Il Mercante che ha venduto informazioni una volta e
non vuole farlo di nuovo.

\secsep

\section{Connessioni Passionali}

All'inizio della sessione, dopo aver scelto i Ruoli, ogni giocatore
dichiara una \textbf{Connessione Passionale}: una relazione specifica
con un altro PG (dello stesso tavolo o, se il GM coordinatore lo
permette, di un altro tavolo) che alimenterà il dramma.

\begin{itemize}[nosep]
  \item La Connessione può essere d'amore, di odio, di debito,
    di sangue, di segreto condiviso
  \item Ogni volta che la Connessione è rilevante in una scena,
    entrambi i personaggi coinvolti guadagnano 1 Passione
  \item Se la Connessione viene spezzata durante la sessione
    (tradimento, morte, separazione definitiva), entrambi guadagnano
    3 Passione contemporaneamente
\end{itemize}

\nota[Suggerimento — Connessioni pre-costruite]{
  Per accelerare la fase di setup, puoi assegnare Connessioni
  pre-costruite basate sui Legami Iniziali dei playbook. Il Cavaliere
  e il Figlio sono già in una relazione maestro-allievo perfetta per
  la Figura Paterna/Materna + il Protagonista. Il Cospiratore e
  la Spia sono già un Rivale naturale.
}

\nota[Prop opzionale — Il nastro della Connessione]{
  Dal vivo, ogni Connessione può essere rappresentata da un nastro o
  cordoncino che lega fisicamente i due giocatori (vedi «Componenti
  Fisici Opzionali» nel capitolo delle regole). Il colore può codificare
  la natura del legame — amore, odio, debito, sangue, segreto — ma anche
  un nastro neutro funziona. Quando la Connessione si spezza, il nastro
  viene reciso: il gesto vale la regola.
}

\subsection{Legare e Sciogliere — l'Azione del Clero}

La Chiesa rivendica il potere di \textit{ligare et solvere}: legare e
sciogliere, in cielo e in terra. Al tavolo del Clero, questo non è solo
teologia.

\regola[Legare e Sciogliere]{
  Un personaggio della fazione \textbf{Clero} con un Ruolo Passionale
  può, \textbf{una volta per atto}, spendere Passione per intervenire
  sulle Connessioni \textit{altrui} — mai sulle proprie:
  \begin{itemize}[nosep]
    \item \textbf{Legare} (spendi 2 Passione) — dichiari una nuova
      Connessione fra altri due personaggi (PG o PNG significativi).
      I due legati non possono rifiutare il legame, ma scelgono loro
      la sua natura. Si tende un nastro fra loro.
    \item \textbf{Sciogliere} (spendi 3 Passione) — spezzi una
      Connessione esistente fra altri due personaggi. Vale come
      Connessione spezzata: entrambi i legati guadagnano 3 Passione.
      Il nastro viene reciso.
  \end{itemize}
  In entrambi i casi il Clero deve giustificare il gesto nella finzione:
  una benedizione, una confessione ascoltata, un matrimonio, una
  scomunica, una verità rivelata. Senza la scena, non c'è il potere.
}

\nota[Per il GM — equilibrio]{
  Legare e Sciogliere è potente: dà al Clero una leva sul dramma di tutti
  gli altri tavoli. È voluto — la Chiesa del 1282 \textit{era} questo. Ma
  se un giocatore del Clero ne abusa, ricordagli che ogni legame imposto
  crea risentimento: i personaggi legati o sciolti contro la loro volontà
  hanno ottime ragioni per farla pagare al Clero, in scena e nel punteggio
  di fazione.
}

\secsep

\section{La Scena della Passione}

Una volta per sessione — tipicamente intorno al secondo Momento di
Congiunzione (La Crisi) — il GM può dichiarare una
\textbf{Scena della Passione}: una breve scena fuori dal flusso
principale dell'azione in cui due personaggi (con Connessione
Passionale tra loro) si trovano soli per 5--10 minuti di gioco.

Nessuna meccanica. Nessun tiro. Solo la scena.

Alla fine della Scena della Passione, entrambi i giocatori dichiarano
cosa è cambiato nel loro personaggio. Possono guadagnare o spendere
Passione liberamente. Il GM prende nota e integra le conseguenze
nel finale.

\nota[Nota di design — Perché questo funziona]{
  Il layer Passionale non rallenta il gioco politico — lo radica.
  Quando i giocatori hanno investimento emotivo nelle relazioni
  personali, le scelte politiche acquistano peso. Il Barone che
  deve scegliere tra la causa e il figlio è più interessante del
  Barone che sceglie tra due opzioni strategiche. Il dramma personale
  \textit{è} la strategia, in Vespri 1282 — perché era così anche
  nella storia vera.
}

\filettodoppio
